\scanpage{119}
\begin{center}- 90 -\end{center}
\margin{n\textsuperscript{o} 371}

tout $i$. Montrons que l'on a $(u(x))^2=0$ pour tout $x\in M$; on peut en effet écrire $x=\sum_{i=1}^n x_i$, où $x_i\in M_i$, donc $(u(x))^2=\sum_{i=1}^n(u(x_i))^2+\sum_{i<j}(u(x_i)u(x_j)+u(x_j)u(x_i))$.

Or, comme $\psi_i(x_i)$ est un élément de degré 1 dans $\bigwedge(M_i)$, on a, pour $i\ne j$, $h_i(\psi_i(x_i))h_j(\psi_j(x_j))=-h_j(\psi_j(x_j))h_i(\psi_i(x_i))$ (§ 3, n\textsuperscript{o} 3, prop.2), d'où $u(x_i)u(x_j)+u(x_j)u(x_i)=0$; par ailleurs $(u(x_i))^2=h_i((\psi_i(x_i))^2)=0$ par définition de $\bigwedge(M_i)$. En vertu de la prop.3 du n\textsuperscript{o} 2, il existe donc un homomorphisme unique d'algèbres graduées $h:\bigwedge(M)\longrightarrow{}^{\unsure{\mathrm g|\mathrm s}}\!\bigotimes_i(\bigwedge(M_i))$ tel que $u=h\circ\psi$ (où $\psi:M\longrightarrow\bigwedge(M)$ est l'injection canonique). Pour $x_i\in M_i$, on a $h(g(u_i(x_i)))=h(g(h_i(\psi_i(x_i))))=h(f_i(\psi_i(x_i)))=h(\psi(j_i(x_i)))=u_i(x_i)$, et comme les $u_i(x_i)$ engendrent l'algèbre ${}^{\unsure{\mathrm g|\mathrm s}}\!\bigotimes_i(\bigwedge(M_i))$, $h\circ g$ est l'identité; de même $g(h(\psi(j_i(x_i))))=g(u_i(x_i))=g(h_i(\psi_i(x_i)))=f_i(\psi_i(x_i))=\psi(j_i(x_i))$, et les éléments $\psi(j_i(x_i))$ engendrent l'algèbre $\bigwedge(M)$, donc $g\circ h$ est l'identité.

COROLLAIRE 1.- \uline{Sous les hypothèses de la prop.7}, \uline{pour tout entier} $m\geq0$, \uline{la puissance extérieure m-ème} $\mathop{\bigwedge}\limits^m M$ \uline{s'identifie canoniquement par} $g$ \uline{à la somme directe des modules} $\bigotimes_{1\leq i\leq n}(\mathop{\bigwedge}\limits^{m_i}M_i)$, \uline{où} $(m_i)_{1\leq i\leq n}$ \uline{parcourt l'ensemble des suites d'entiers} $\geq0$ \uline{telles que} $\sum_{i=1}^n m_i=m$.

Cet isomorphisme s'explicite de la façon suivante: considérons $m$ éléments $x_j=\sum_{i=1}^n x_{ji}$ $(1\leq j\leq m)$ de $M$, avec $x_{ji}\in M_i$ pour $1\leq i\leq n$, et pour toute application $\sigma$ de $[\![1,m]\!]$ dans $[\![1,n]\!]$, formons le m-vecteur
\[
\text{(7)}\qquad x_{1,\sigma(1)}\wedge x_{2,\sigma(2)}\wedge\cdots\wedge x_{m,\sigma(m)}\ .
\]
La donnée de l'application $\sigma$ équivaut à la donnée, pour chaque $i\in[\![1,n]\!]$, d'une partie $H_i=\sigma^{-1}(i)$ de $[\![1,m]\!]$, celles des $H_i$ non vides formant une \uline{partition} de $[\![1,m]\!]$; si $H_i$ est vide, posons $x_{H_i,i}=1$, et dans le cas contraire, $x_{H_i,i}=x_{j_1,i}\wedge x_{j_2,i}\wedge\cdots\wedge x_{j_{r_i},i}$, où $(j_k)_{1\leq k\leq r_i}$ est la suite
\margin{118}
